\documentclass[11pt]{article}
\usepackage[T1]{fontenc}
\usepackage{textcomp}
\usepackage[margin=0.75in]{geometry}
\usepackage{amsmath,amssymb,amsthm}
\usepackage{array,booktabs}
\usepackage{hyperref}
\setlength{\parindent}{0pt}
\newtheorem{theorem}{Theorem}
\theoremstyle{definition}
\newtheorem{definition}{Definition}
\title{Flow Proof Artifact: prop-13-no-double-touch-internal-external}
\author{Generated by \texttt{flow doc proof}}
\date{Flow Proof Book}
\begin{document}
\maketitle
A circle does not touch another internally or externally at more than one point.\\[0.5em]
\textbf{Source.} euclid — Elements, Book III, Proposition 13\\[0.5em]
\bigskip
\noindent{\large \textbf{Derived fact 1.} \textit{A circle does not touch another internally or externally at more than one point} \hfill the Euclidean plane \cdot Euclid Book III \cdot Proposition 13: circles cannot touch at more than one point internally or externally}\par
\smallskip\noindent\textit{Coordinate: the Euclidean plane \cdot Euclid Book III \cdot Proposition 13: circles cannot touch at more than one point internally or externally}\par
\smallskip\noindent\textit{Source: euclid — Elements, Book III, Proposition 13}\par
\smallskip\noindent\textit{Built on: proposition 10: one circle cannot touch another at more than one point, for Euclid Book III on the Euclidean plane, proposition 11: internally touching circles have collinear centers through the contact, for Euclid Book III on the Euclidean plane}\par
\medskip
\noindent\textbf{Goal.} A circle does not touch another internally or externally at more than one point\par
\noindent$\displaystyle \text{no double touch}$\par
\medskip
\noindent
\begin{tabular}{@{} >{\raggedright\arraybackslash}p{0.06\textwidth} >{\raggedright\arraybackslash}p{0.47\textwidth} >{\raggedleft\arraybackslash}p{0.41\textwidth} @{}}
\textbf{\#} & \textbf{Proof} & \textbf{Mathematics} \\
\midrule
\textbf{1.} & We prove that proposition 13: circles cannot touch at more than one point internally or externally for Euclid Book III the Euclidean plane. & \\[0.45em]
\textbf{2.} & Let C1 = first\_circle. & \\[0.45em]
\textbf{3.} & Let C2 = second\_circle. & \\[0.45em]
\textbf{4.} & From step 2 and step 3, we can deduce that at most one touch point. & $\displaystyle \text{at most one touch point}$ \\[0.45em]
\textbf{5.} & We invoke the derived fact governing Euclid Book III the Euclidean plane: proposition 10: one circle cannot touch another at more than one point, for Euclid Book III on the Euclidean plane. & \\[0.45em]
\textbf{6.} & We invoke the derived fact governing Euclid Book III the Euclidean plane: proposition 11: internally touching circles have collinear centers through the contact, for Euclid Book III on the Euclidean plane. & \\[0.45em]
\textbf{7.} & From step 5 and step 6, this implies no double touch. Hence proven. & $\displaystyle \text{no double touch}$ \\[0.45em]
\bottomrule
\end{tabular}
\label{thm:Geometry-Euclid Book III-proposition_13_circles_cannot_touch_at_more_than_one_point_internally_or_externally}
\par\medskip\hrule
\end{document}
